% =====================================================================
%  CV — Huawei Canada, Noah's Ark Laboratory, Markham
%        Co-op Researcher — Physical AI (embodied / spatial AI).
%  Dérivé de master-cv.md (source de vérité unique)
%
%  CIBLAGE — le poste (embodied / spatial AI) demande : Deep Learning
%  (training, regularization, generalization), Python + PyTorch, et cite
%  robotique et LLMs COMME ATOUTS, pas exigences. L'angle honnête :
%   - Le MOTEUR DE DYNAMIQUE RIGIDE DIFFÉRENTIABLE et la SORTIE EN FORCES
%     sous cône de friction = du contrôle bas niveau physiquement admissible.
%     C'est le point de contact réel avec le Physical AI ; il passe en tête.
%   - « world / spatial model » ↔ le corpus render-then-estimate et la
%     profondeur monoculaire = raisonnement spatial à partir de pixels.
%   - PyTorch + DL components : en tête des compétences, ce sont leurs mots.
%
%  CE QUI N'EST PAS REVENDIQUÉ, faute de preuve : robotique, LLMs, RL,
%  world models, sim-to-real. Le poste les cite comme atouts ; les inventer
%  se paierait au premier entretien. La proximité est nommée honnêtement
%  (contrôle physique, dynamique), la possession ne l'est pas.
%
%  INTERDITS (master-cv.md §11) : projet clical, −58 %, PFA 82,3 %,
%  audience Medium, Google Final Round, projets perso, noms de clients,
%  rejet antérieur du papier ICLR, canal Talents Handicap, mandarin
%  « working knowledge » (niveau réel : débutant).
%
%  ⚠ AVANT ENVOI — deux points à trancher, voir la conversation :
%    (1) ces co-ops sont annoncés 8 à 16 mois ; la disponibilité réelle est
%        février → octobre 2027, soit 8 mois. Le formulaire ne doit cocher
%        que « 8 months », jamais 12 ni 16.
%    (2) autorisation de travail au Canada pour 2027 : rien n'est écrit ici
%        tant que le statut n'est pas établi.
% =====================================================================

\documentclass[10pt, letterpaper]{article}

\usepackage[
    ignoreheadfoot,
    top=0.7 cm,
    bottom=0.55 cm,
    left=1.15 cm,
    right=1.15 cm,
    footskip=0.8 cm,
]{geometry}
\usepackage{titlesec}
\usepackage{enumitem}
\usepackage{fontawesome5}
\usepackage[dvipsnames]{xcolor}
\usepackage{textcomp}

\definecolor{MyRed}{RGB}{160, 20, 40}

\usepackage[
    pdftitle={Melissa Colin — CV — Huawei Canada, Co-op Researcher, Physical AI},
    pdfauthor={Melissa Colin},
    colorlinks=true,
    urlcolor=MyRed
]{hyperref}

\usepackage{charter}

\pagestyle{empty}
\setlength{\parindent}{0pt}

\titleformat{\section}
  {\normalsize\bfseries\scshape\raggedright\color{MyRed}}{}{0em}{}[\titlerule]
\titlespacing{\section}{0pt}{6pt}{4pt}

\setlist[itemize]{leftmargin=*, noitemsep, topsep=1pt, partopsep=0pt, parsep=0pt}

\newcommand{\entry}[4]{%
  \noindent \textbf{#1} \hfill #2 \\
  \textit{#3} \hfill \textit{#4} \par
}

\IfFileExists{personal.tex}{\input{personal.tex}}{}
\providecommand{\myphone}{+33 7 82 52 XX XX}

\begin{document}

% ---------------------------------------------------------------- HEADER
\begin{center}
    {\fontsize{19}{19}\selectfont \textbf{Mélissa COLIN}} \\[3pt]
    {\normalsize Deep Learning \& Differentiable Physics $\cdot$ Physically Grounded 3D Human Motion} \\[3pt]

    \small{
    \faMapMarker* Bordeaux, France \;|\;
    \faPhone* \myphone{} \;|\;
    \faEnvelope~\href{mailto:contact@melissacolin.ai}{contact@melissacolin.ai} \;|\;
    \faGlobe~\href{https://melissacolin.ai}{melissacolin.ai} \;|\;
    \faGithub~\href{https://github.com/melissa-colin}{melissa-colin} \;|\;
    \faGraduationCap~\href{https://scholar.google.com/citations?user=7r7iFpsAAAAJ}{Scholar}
    }
\end{center}

\vspace{1pt}

{\small \textit{Final-year MEng student in computer science, first author on a CVPR 2027
submission. I build a differentiable physics engine and neural models that output forces
under physical constraints, so that 3D human motion recovered from video is physically
admissible by construction rather than by penalty. Expected graduation 09/27, available from
February 2027 for a co-op of up to eight months, through October 2027.}}

% ---------------------------------------------------------------- PUBLICATIONS
\section{Publications}
\begin{itemize}
    \item \textbf{Colin, M.}, Wang, Y., Cheng, L. \textit{InterForce: Physically Admissible Multi-Person Motion Refinement.} First author, in preparation, targeting \textbf{CVPR 2027}.
    \item \textbf{Colin, M.}, Chraibi Kaadoud, I. \textit{Performances and Explainability of ViT and CNN Architectures: An Empirical Study Using LIME, SHAP and Grad-CAM.} RJCIA / PFIA 2024, La Rochelle. First author, published, presented orally. \href{https://hal.science/hal-04641791}{[HAL]}
    \item \textbf{Co-author}, 3D group dance generation. In preparation, targeting ICLR 2027.
\end{itemize}

% ---------------------------------------------------------------- RESEARCH
\section{Research Experience}

\entry{Vision and Learning Lab, University of Alberta (Prof.\ Li Cheng)}{May 2026 -- Sep 2026}%
{Visiting Research Student, \textbf{Mitacs Globalink Research Award}}{Edmonton, Canada}
\begin{itemize}
    \item Wrote a \textbf{differentiable rigid-body dynamics engine from scratch} --- recursive Newton--Euler inverse dynamics, composite-rigid-body mass matrix, implicit actuation, joint limits --- and \textbf{backpropagate through it inside the training loop}. Checked against nine analytic invariants to machine precision.
    \item Designed a model that \textbf{outputs forces, not poses}, projected into their friction cone, so a physically impossible motion becomes \textbf{unrepresentable rather than penalised}: this is low-level control constrained to be physically admissible. Clips a simulator replays \textbf{without the character falling: 19\% $\to$ 79\%}.
    \item Designed a \textbf{permutation-equivariant transformer} whose attention is factorised over people, joints and time, so a multi-person scene is handled without fixing the number of people.
    \item Built the \textbf{27{,}000-clip corpus} by render-then-estimate and recovered metric-scale spatial structure (monocular depth, occlusion ordering) straight from pixels: \textbf{$-$33\% inter-person penetration} over the whole corpus, purely algorithmic.
\end{itemize}

\vspace{1pt}

\entry{ENSEIRB-MATMECA --- Research Initiation, 17/20}{Feb -- May 2026}%
{Reproduction and critical study of \textit{InterMask} (ICLR 2025)}{Bordeaux, France}
\begin{itemize}
    \item Reproduced the \textbf{VQ-VAE} baseline (FID 1.034 $\to$ \textbf{0.918}), then tested two hypotheses on the physical constraints of the loss. Neither improved FID or R-Precision; \textbf{reported the negative result} instead of dropping it.
\end{itemize}

% ---------------------------------------------------------------- INDUSTRY
\section{Industry Experience}

\entry{Sector Group}{Jun -- Aug 2025}{R\&D Engineer, AI and Safety --- information extraction from \textbf{PDFs with no text layer}, which ruled out parsing and forced vision and document processing together}{Paris, France}

\vspace{1pt}

\entry{Cali Intelligences (deeptech start-up, real-time video analytics)}{Jan 2023 -- Sep 2024}%
{Applied AI Intern, then AI Developer (apprenticeship)}{Talence, France}
\begin{itemize}
    \item Took detection accuracy from \textbf{60\% to 90\%} by designing and building the entire training pipeline (YOLOv7 $\to$ YOLOv8): ingestion, label conversion, an image augmentation library written from scratch in OpenCV, pre-training validation, and YAML-driven orchestration.
    \item Shipped it to production on Docker, GitHub Actions and K3s.
\end{itemize}

% ---------------------------------------------------------------- EDUCATION
\section{Education}

\entry{ENSEIRB-MATMECA, Bordeaux INP}{Sep 2024 -- Expected 09/27}%
{MEng in Computer Science (\textit{Diplôme d'ingénieur}), AI track, final year}{Bordeaux, France}
\begin{itemize}
    \item Ranked \textbf{5\textsuperscript{th}/81 (top 6\%)} in semester 8 and \textbf{3\textsuperscript{rd}/93} in semester 6 (French /20 scale, transcript available). \textbf{Deep Learning 18/20} $\cdot$ \textbf{Research Initiation 17/20}.
\end{itemize}

\vspace{1pt}

\entry{EPSI Bordeaux}{Sep 2021 -- Sep 2024}%
{BSc in Artificial Intelligence \& Data Science, ranked \textbf{1\textsuperscript{st}/10} in final year}{Bordeaux, France}

% ---------------------------------------------------------------- SKILLS
\section{Technical Skills}
\begin{itemize}
    \item \textbf{Deep learning:} \textbf{PyTorch} (advanced) --- training, regularization, generalization $\cdot$ Transformers and ViT $\cdot$ VQ-VAE $\cdot$ CNNs $\cdot$ Keras/TensorFlow $\cdot$ LoRA/PEFT $\cdot$ written from scratch: backpropagation, PCA, K-Means, logistic regression.
    \item \textbf{Physics \& control:} \textbf{differentiable rigid-body dynamics} (RNEA, mass matrix, joint limits) $\cdot$ force-level control under friction constraints $\cdot$ Isaac Gym rigid-body simulation $\cdot$ MCTS / self-play.
    \item \textbf{3D vision:} SMPL / SMPL-X, forward kinematics, axis-angle and 6D rotations $\cdot$ monocular depth $\cdot$ camera calibration, occlusion ordering $\cdot$ OpenCV.
    \item \textbf{Programming \& foundations:} \textbf{Python} $\cdot$ C $\cdot$ C++ $\cdot$ Java $\cdot$ JavaScript $\cdot$ SQL $\cdot$ Bash $\cdot$ Linux/POSIX, Git, Docker, multi-GPU clusters $\cdot$ numerical linear algebra, probability, optimisation, NP-completeness.
    \item \textbf{Languages:} French (native) $\cdot$ English \textbf{B2, TOEIC 840/990}, placement conducted entirely in English $\cdot$ Mandarin (elementary).
\end{itemize}

\end{document}
